\chapter{PENDAHULUAN}
\label{chap:pendahuluan}

% Ubah bagian-bagian berikut dengan isi dari pendahuluan

Penelitian ini dilatarbelakangi oleh \lipsum[1][1-5]

\section{Latar Belakang}
\label{sec:latarbelakang}

Penyakit yang berkontribusi sebagai penyebab kematian terbesar di indonesia diantaranya adalah penyakit stroke dengan pesentase 19.42 persen dan jantung iskemik dengan persentase 14.38 persen. tidak hanya di indonesia, kedua penyakit kardiovaskular ini juga menjadi per- hatian dari lembaga kesehatan dunia karena penyakit stroke berkontribusi sebesar 16.17 persen total kematian dunia dan penyakit jantung iskemik berkontribusi sebesar 11.59 persen total kematian dunia [1].

penyakit stroke adalah penyakit yang memengaruhi arteri yang menuju dan/atau di dalam otak. Stroke terjadi ketika pembuluh darah yang membawa oksigen dan nutrisi ke otak terhalang oleh gumpalan darah atau pecah. Ketika hal itu terjadi, sebagian dari otak tidak bisa menda- patkan darah dan oksigen yang dibutuhkannya, sehingga sel-sel dalam otak dan/atau otaknya mati [2]. Sedangkan penyakit jantung iskemik adalah kondisi medis yang terjadi ketika aliran darah ke jantung berkurang akibat penyumbatan sebagian atau seluruh arteri koroner. kedua penyakit ini memiliki penyebab yang sama yaitu sumbatan pada pembuluh darah [3].

Teknik yang umum dilakukan untuk mengecek kondisi penyumbatan pada pembuluh darah adalah Diagnostic coronary angiography, teknik ini juga bisa untuk menilai seberapa parah sum- batan pada pembuluh darah untuk menentukan langkah selnjutnya. Diagnostic coronary angiog- raphy dilakukan dengan memberikan anestasi lokal pada daerah tubuh yang akan dimasukkan kateter, biasanya lengan atau paha. dari area itu kateter yang berupa tabung tipis dan fleksibel akan dimasukkan ke pembuluh darah dan diarahkan menuju jantung. Sinar x-ray digunakan un- tuk membantu mengarahkan kateter, setelah kateter terpasang akan disuntikkan cairan pewarna yang disebut cairan kontras. gambar sinar x-ray akan memperlihatkan bagaimana cairan terse- but bergerak di dalam pembuluh darah, dan membantu menemukan sumbatan pada pembuluh darah. Teknik ini dapat menimbulkan rasa tidak nyaman pada pasien dan resiko terpapar radiasi dari x-ray.

Penerapan realitas virtual (VR) dalam dunia pendidikan dan medis memiliki potensi yang sangat besar untuk merevolusi cara belajar dan memberikan layanan kesehatan. Dalam bidang medis, VR memungkinkan pelatihan yang lebih mendalam dan realistis bagi tenaga medis, termasuk simulasi operasi yang aman dan efektif tanpa risiko terhadap pasien. Teknologi ini memungkinkan para profesional medis untuk berlatih prosedur yang kompleks dalam lingkun- gan yang dikendalikan, sehingga meningkatkan keterampilan dan kepercayaan diri mereka. Se- lain itu, VR digunakan dalam rehabilitasi pasien, memungkinkan mereka untuk berlatih gerakan dan aktivitas dalam lingkungan virtual yang dapat disesuaikan dengan kebutuhan individu.

Dalam konteks pendidikan medis, VR juga memungkinkan pengalaman belajar yang lebih imersif dan interaktif. Mahasiswa kedokteran dapat memanfaatkan teknologi ini untuk mem- pelajari anatomi manusia secara lebih mendetail melalui simulasi 3D, serta berpartisipasi dalam skenario klinis yang realistis. Dengan adanya pandemi Covid-19, adaptasi dan inovasi dalam penggunaan teknologi VR semakin meningkat, mendukung teleworking dan pelatihan jarak jauh. Pelatihan yang efektif dalam penggunaan teknologi ini sangat penting untuk memastikan manfaat maksimal dalam pendidikan medis. Di masa depan, VR diharapkan menjadi salah satu metode utama dalam proses pembelajaran dan pelatihan medis, terutama untuk pembelajaran mandiri dan di area di mana teknologi VR lebih berkembang [4].


\lipsum[2]

\section{Permasalahan}
\label{sec:permasalahan}

Dalam dunia kesehatan, penyakit yang disebabkan oleh penyumbatan pada pembuluh darah merupakan salah satu kontibutor kematian terbesar dengan persentase lebih dari 25 persen total kematian dunia.

Kesadaran masyarakat tentang pentingnya menjaga kesehatan pembuluh darah dinilai ku- rang, padahal kesehatan pembuluh darah sangat penting karena berperan dalam mengatur aliran darah yang membawa oksigen dan nutrisi ke seluruh tubuh. Sistem pemeriksaan yang ada meli- batkan kateter yang dimasukkan ke dalam tubuh dan juga sinar x-ray yang bisa memberikan rasa tidak nyaman pada pasien serta kemungkinan resiko radiasi dari x-ray.

Penggunaan VR dalam dunia edukasi dapat bermanfaat dengan mensimulasikan fisiologis dalam tubuh manusia khususnya peredaran darah manusia.


\section{Tujuan}
\label{sec:Tujuan}

Tujuan utama dalam penelitian ini adalah mengembangkan aplikasi untuk menvisualisas- ikan tekanan pada pembuluh darah dengan menggunakan virtual reality sebagai media visu- alisasi. Dengan demikian, penelitian ini bertujuan untuk memungkinkan pengguna tahu apa yang sedang terjadi di dalam pembuluh darah ketika kondisi normal dan ketika kondisi dengan sumbatan. Aplikasi ini diharapkan dapat digunakan sebagai media edukasi dan promosi tentang pentingnya menjaga kesehatan pembuluh darah kepada masyarakat.


% \begin{enumerate}[nolistsep]

%   \item Membuat \lipsum[2][1-3]

%   \item \lipsum[3][1-3]

% \end{enumerate}

\section{Batasan Masalah}
\label{sec:batasanmasalah}

Ruang lingkup dalam penelitian ini diantaranya adalah:

\begin{enumerate}[nolistsep]

      \item Visualisasi yang akan dilakukan hanya pada arteri dan vena tanpa ada kapiler.
      \item Visualisasi memiliki 5 tingkatan, yaitu tinggi, cukup tinggi, menengah, cukup rendah, dan rendah, bukan dalam bentuk angka.
      \item Sumbatan dalam pembuluh darah memiliki 3 tingkatan, yaitu kecil, sedang, dan besar.
      \item Visualisasi \emph{virtual reality} dilakukan di perangkat Oculus Meta 2 dengan memanfaatkan fitur interaktif yaitu \emph{hand tracking}.

\end{enumerate}

\section{Sistematika Penulisan}
\label{sec:sistematikapenulisan}

Laporan penelitian tugas akhir ini terbagi menjadi \lipsum[1][1-3] yaitu:

\begin{enumerate}[nolistsep]

  \item \textbf{BAB I Pendahuluan}

        Bab ini berisi \lipsum[2][1-5]

        \vspace{2ex}

  \item \textbf{BAB II Tinjauan Pustaka}

        Bab ini berisi \lipsum[3][1-5]

        \vspace{2ex}

  \item \textbf{BAB III Desain dan Implementasi Sistem}

        Bab ini berisi \lipsum[4][1-5]

        \vspace{2ex}

  \item \textbf{BAB IV Pengujian dan Analisa}

        Bab ini berisi \lipsum[5][1-5]

        \vspace{2ex}

  \item \textbf{BAB V Penutup}

        Bab ini berisi \lipsum[6][1-5]

\end{enumerate}
